\chapter{生成的本质：从图像到分布}
\label{ch:concepts}

\readingguide{本章先建立一个能够反复使用的直觉：图像生成要学会的，不是一张代表性图片，而是产生许多合理图片的规律。我们会依次讨论分布、高维几何、加噪与去噪、隐空间和条件生成，再把扩散放回生成模型的全景中。读完后，你应当能够解释为什么一个预测噪声的网络可以参与图像生成，也能识别这种解释的边界。涉及的概率术语可点击跳转到附录；完整推导将在第\ref{ch:math}章展开。}

\section{先问一个看似简单的问题：什么叫会生成？}
\label{sec:c-generation}

假设我们给模型一百万张猫的照片，然后要求它生成一张新照片。最容易想到的方案是输出这批照片的平均值。这个方案在某种意义下十分合理：如果使用像素平方误差，且模型只能输出一张固定图片，平均图像就是最优答案。但它往往是一团模糊的轮廓。猫的头可以向左，也可以向右；眼睛可以睁开，也可以闭上；背景可以是窗台，也可以是草地。这些可能性逐像素平均之后，不再构成其中任何一个合理场景。

问题出在任务的描述。我们希望模型生成一个可能发生的结果，而不是在所有可能结果之间折中。一张照片只能呈现一个视角、一种姿态和一种光照；一套生成规则却应当容纳许多视角、姿态和光照。描述这套规则最自然的数学对象，就是\term{概率分布}{gl:distribution}。

用 $X$ 表示随机图像，用 $x$ 表示它的一次具体取值。即使我们最终只想得到一张图片，也应当把目标写成
\[
  \widehat X\sim p_\theta,\qquad p_\theta\approx \pdata,
\]
其中 $\pdata$ 是我们想要建模的数据分布，$p_\theta$ 是参数为 $\theta$ 的模型产生的分布。符号 $\sim$ 读作“服从”；它把一次生成和背后的重复试验联系起来。我们真正关心的是：如果多次运行模型，出现的图片是否合理，是否有足够的变化，各类内容出现的比例是否符合目标。

\begin{examplebox}[unbreakable]{平均答案与随机答案}
一个数据集只有两个标量值 $-2$ 和 $2$，各占一半。输出 $0$ 的模型实现了最小的常数平方误差，却从不输出数据中可能出现的值。另一个模型先掷一枚公平硬币，再输出 $-2$ 或 $2$，才正确复现了这个分布。图像中的“不同姿势”与这个例子的两个端点相似，只是维数高得多。
\end{examplebox}

由此可以区分三个经常混在一起的目标。重建要求输出尽量接近某张给定图片；预测要求根据条件判断未来或缺失的内容；生成要求从一组可能性中抽取一个结果。它们可以共享网络，但评价标准不同。重建的答案经常可以与原图逐像素比较，无条件生成的答案没有这样一个唯一参照。把生成当作单一答案回归，模糊便成为一种可以预期的结果。

“复现分布”也不意味着逐张复制训练集。如果训练集为 $\{x^{(1)},\ldots,x^{(n)}\}$，它对应的经验分布给每张训练图分配质量 $1/n$。随机抽取并原样输出一张训练图，可以完美复现这个经验分布，却没有学会在未见情形下生成。实际学习希望利用有限样本、模型结构和训练目标，推断更丰富的总体规律。这就是生成模型里的泛化问题。训练样本本身有偏，学到的规律也可能偏；网络容量足够大时，它同样可能记忆样本。扩散并没有天然免除这些问题。

还应当明确“真实”的含义。互联网上收集的图片分布不是自然界全部可能图像的分布，更不是人类审美的统一标准。相机、网页选择、标注方式和筛选过程都在改变它。若数据集中人物通常居中，模型就可能把“居中”当作高概率规律。我们在全书使用 $\pdata$，指的是指定任务与数据获取过程下的目标分布；它不是一个脱离任务而客观存在的完美图库。

\begin{intuitionbox}{生成能力是重复运行的性质}
单张图片可以漂亮，但一个模型是否学得好，要看它产生的整批结果。只输出一张漂亮图片的模型，可能缺少覆盖；能产生很多不同图片的模型，也可能经常输出不合理内容。合理性、覆盖范围和概率分配是彼此相关、又不能完全互相替代的要求。
\end{intuitionbox}

\section{分布、密度和采样：三个不同的问题}
\label{sec:c-distribution}

\subsection{分布描述概率，密度描述概率如何铺开}

一张 $H\times W$ 的 RGB 图像可表示为 $d=3HW$ 维向量。对于八位像素，严格说来它属于一个巨大的离散集合；若把像素归一化并作为实数处理，便常写 $x\in\R^d$。这两种建模方式有不同的概率含义，不能在计算\term{似然}{gl:likelihood}时随意交换。连续模型处理离散图像时，还需要对量化区间或去量化作出约定。

分布 $P$ 为事件分配概率，例如“图像里有两只猫”或“某个像素位于某个区间”。当分布有\term{概率密度}{gl:density} $p(x)$ 时，
\[
  \mathbb P(X\in A)=\int_A p(x)\,\dd x.
\]
因此 $p(x)$ 本身不是“恰好生成图片 $x$ 的概率”。对通常的连续分布，单点概率为零；一个小区域的概率，才近似等于密度乘以区域体积。密度也可以大于一，因为它有“每单位体积”的含义。只要积分为一，就是合法的密度。

这一区分在高维空间尤其关键。密度比较的是单位体积中的集中程度，采样频率比较的是整个区域累积的概率质量。面积很小的高密度区域，可能比面积巨大的中等密度区域更少被抽到。若把生成理解为“寻找 $p(x)$ 最大的图片”，就会丢掉这种体积效应，也会把多样性缩成一个优化问题。

\begin{examplebox}{最密的点，未必是最典型的样本}
设 $Z\sim\N(0,I_d)$。密度在原点最大，但
\[
  \E\norm{Z}^2=d,\qquad
  \operatorname{Var}(\norm{Z}^2)=2d.
\]
由切比雪夫不等式，对任意 $\delta>0$，
\[
 \mathbb P\!\left(\left|\frac{\norm{Z}^2}{d}-1\right|>\delta\right)
 \leq \frac{2}{d\delta^2}.
\]
维数很大时，大部分样本的半径接近 $\sqrt d$，而不是接近零。原点密度高，原点附近足够小的球却几乎没有总体概率质量。
\end{examplebox}

这个例子不只是一个概率练习。扩散模型常从高维高斯噪声出发；所谓“随机噪声的常见样子”，就是许多坐标同时有中等幅度的随机波动，而不是所有坐标接近零。类似地，一张图片在模型下具有较高密度，未必就具有更好的视觉质量。密度还依赖坐标和量化约定，视觉质量则依赖结构、语义和感知。它们需要联系起来研究，不能当作同一个指标。

\subsection{知道密度和会采样并不等价}

假设有人给你函数 $p(x)$，你能计算任意图片的密度。你仍然不一定能高效地产生它的样本。高维空间里随机乱找，几乎不会碰到合理图片；寻找最大值，又只是在做优化。反过来，一个网络可以把随机向量直接变成图片，因而很容易采样，但未必能给这张图片计算归一化密度。生成模型的不同家族，往往就在这两个能力之间作不同安排。

对确定性生成器，可以写
\[
 Z\sim\pnoise,\qquad \widehat X=G_\theta(Z).
\]
生成器把输入分布“推送”为输出分布。这里的数学含义是
\[
 \mathbb P(\widehat X\in A)
 =\mathbb P\!\left(Z\in G_\theta^{-1}(A)\right).
\]
这里 $G_\theta^{-1}(A)$ 表示所有会被映射进 $A$ 的输入集合，称为原像；它不要求 $G_\theta$ 存在逆函数。即使 $G_\theta$ 的每一步完全确定，只要 $Z$ 是随机的，输出仍然是随机变量。一个固定函数可以产生无限多的结果，并不需要为每个结果单独存储一个模板。

这一观察能提前解开一个常见困惑：扩散采样一定要在每一步加入随机噪声吗？答案是否定的。某些采样方式只在初始状态随机，随后沿确定性轨迹运动；另一些方式在中途继续注入随机性。只要最终的分布正确，两种方式都可以进行生成。第\ref{ch:math}章将解释随机微分方程与概率流常微分方程为什么能在适当条件下拥有相同的时刻分布。

分布还可通过样本来认识，而不必先知道一个可计算的公式。训练集提供的就是这种信息：我们会从 $\pdata$ 取样，却不知道 $\pdata(x)$ 的解析表达式。有效的训练目标必须适应这个现实。扩散之所以吸引人，一个重要原因是：它能利用已知的加噪规则，把“只有干净样本”转化成大量有监督的预测任务。

\subsection{边缘、条件和路径分别在说什么}

假设 $X_t$ 是某个时间的含噪图像。$p_t(x)$ 描述这一时刻的一整群样本，这是时刻\term{边缘分布}{gl:distribution}。$q(x_t\mid x_0)$ 描述给定一张原图之后，加噪的结果怎样变化，这是\term{条件分布}{gl:conditional}。$q(x_0,x_1,\ldots,x_K)$ 则描述一条完整随机轨迹的联合规律。边缘、条件和联合分布不应混为一谈。

仅知道每一个时刻的边缘分布，不会唯一决定样本之间怎样连线。想象两张纸上各有一百个点，第一张代表初始噪声，第二张代表最终数据。你可以按不同规则把两张纸上的点配对，最终点云完全相同，运动路线却不同。第\ref{ch:methods}章中不同的采样过程和第\ref{sec:flowmatching}节中的路径设计，都将使用这一区分。

\section{为什么图像分布这么难学？}
\label{sec:c-geometry}

\subsection{维数高，更困难的是结构强}

一张 $256\times256$ 的 RGB 图像就有 $196\,608$ 个坐标。若每个坐标独立取一个像素值，所得图片几乎总是杂乱的噪声。自然图像的像素之间存在大量依赖：物体边缘连续，光照在空间上变化有规律，人的左右肢体受共同姿态约束，文字必须由一致的笔画组成。困难不只是坐标多，而是需要把这些约束同时满足。

“局部看起来合理”也不足以保证全局合理。每一小块都像皮毛，整张图却可能有三条尾巴；每个字母都有清楚的线条，一行文字却可能没有拼成任何词。反过来，全局构图正确，细部纹理也可能不可信。有效模型需要在不同尺度上协调信息。这为后文的 U-Net、多尺度特征以及注意力机制提供了问题背景，但概率原理并不规定必须使用某一种网络。

有限数据给学习增加了另一层困难。设想给每个坐标划十个格子，再统计每个格子组合的出现频率，便有 $10^d$ 种组合；现实数据远远无法覆盖。神经网络并非消除了这个指数级问题，而是用共享参数和结构偏好压缩了需要学习的规律。某种边缘检测模式可以在许多位置复用；某种“物体与阴影相协调”的关系，可以跨越许多具体图片迁移。

\subsection{流形是一幅有用的地图，但不是字面的全部真相}

一个常见想法是，自然图像集中在高维空间中的较低维结构附近。相机位置、物体姿态、光照和材料等因素的变化，可以在像素空间画出连续轨迹。我们用\term{数据流形}{gl:manifold}作为理解这种几何结构的语言：周围空间很大，而有意义的变化只占其中某些方向。

例如，固定一只手写数字的风格，只连续改变它的倾角，图片会沿一条曲线变化。改变倾角和笔画粗细，可能形成一个局部二维曲面。随机改变每一个像素，则大多沿着离开这个曲面的方向运动。对这个理想例子，去噪可以理解为把偏离曲面的观测拉回附近，同时保留曲面内部的有效变化。

不过，真实图像还包含传感器噪声、压缩误差、遮挡、离散类别和复杂组合。不同类别可能对应不同区域；某些变化会产生不连续的拓扑变化。不能因此断言所有图像恰好位于一张光滑、低维、无边界的流形上。“靠近结构化集合”往往比“严格落在一个流形上”更符合实际。后面有些数学公式需要完整空间中的平滑密度，正是加噪帮助我们获得了这种良好的对象。\cite{r_ncsn}

这里还能区分几何和概率。知道哪里有合理图片，只告诉了我们分布的支持或高概率区域；还要知道各区域应分到多少概率质量。两张同样合理的图片，可以在数据中拥有完全不同的出现频率。把样本投影到一个结构化集合，并不自动实现正确的生成分布。

\begin{figure}[htbp]
\centering
\begin{tikzpicture}[x=0.85cm,y=0.85cm]
  \draw[->,gray] (-0.4,0) -- (6.2,0) node[right] {$x^{(1)}$};
  \draw[->,gray] (0,-0.4) -- (0,3.6) node[above] {$x^{(2)}$};
  \draw[line width=1.4pt,blue!70!black]
    plot[smooth] coordinates {(0.45,0.7) (1.1,1.7) (2,2.15) (3,1.55) (4,1.45) (5.4,2.8)};
  \foreach \px/\py in {0.55/0.86,1.12/1.74,1.8/2.14,2.5/1.94,3.2/1.5,4.25/1.63,5.18/2.59}
    \fill[blue!70!black] (\px,\py) circle (2pt);
  \foreach \px/\py/\qx/\qy in {1.05/0.65/1.3/1.88,2.0/2.98/2.0/2.15,3.2/2.55/3.25/1.5,4.75/0.8/4.45/1.8}
    {\fill[orange!85!black] (\px,\py) circle (2.1pt);
     \draw[->,orange!85!black,line width=0.75pt] (\px,\py) -- (\qx,\qy);}
  \node[align=left,font=\small] at (4.5,3.25) {蓝线：结构化数据区域\\橙点：受到扰动的观测};
  \node[font=\small,align=center] at (3,-0.8) {箭头仅示意局部纠正方向；\\沿蓝线的概率分配仍然需要学习。};
\end{tikzpicture}
\caption{二维图示帮助理解高维几何。真实图像通常更复杂，箭头也不必是欧氏最近点投影。图为原创教学示意。}
\label{fig:c-manifold}
\end{figure}

\subsection{没有一条坐标轴天然叫“猫的姿势”}

像素坐标里，改变姿势会同时改变许多位置的颜色；改动一个像素通常不能完成一个清楚的语义变化。这种坐标与语义之间的不匹配，是高维学习的又一个难点。学习到的特征可以让某些规律更容易表达，但也不保证每个特征维度只控制一个独立概念。

因此我们应当谨慎对待“模型学到了语义空间”这句话。它可能表示模型的内部表示对语义有用，也可能只是表示相似内容在某种特征距离下比较接近。它并不自动意味着坐标可解释、变量互相独立、任意插值都合理。即使是一个效果很好的生成模型，仍可能把姿势、背景和风格纠缠在一起。

\section{加噪：先把复杂问题变成一串较温和的问题}
\label{sec:c-noising}

\subsection{扩散不是把一张图均匀地涂成灰色}

扩散模型的基本安排是：规定一个把数据逐渐扰乱的过程，然后学习怎样反向推进。这个想法在早期扩散概率模型中已经出现；DDPM 则给出了影响深远的训练与采样形式。\cite{c_sohl2015diffusion,r_ddpm}

先不讨论具体时间步，考虑
\begin{equation}
  X_t=a(t)X_0+b(t)\eps,\qquad
  X_0\sim\pdata,\quad \eps\sim\N(0,I_d),
  \label{eq:c-noising}
\end{equation}
其中 $X_0$ 与 $\eps$ 独立。$a(t)$ 控制保留多少原信号，$b(t)$ 控制注入多少噪声。在常见设置里，$t$ 从 $0$ 向 $1$ 增加，信号减弱、噪声增强。本书约定生成过程沿相反方向，从 $t=1$ 推进到 $t=0$。不同方法的系数和路径可以不同，不能只凭这一个写法判定整个过程。

对于固定原图 $x_0$，$X_t$ 的条件分布是
\[
 q_t(x\mid x_0)=\N\!\left(x;\,a(t)x_0,b(t)^2I_d\right).
\]
这句话很容易理解：每张原图周围都有一团高斯云。对所有原图取平均，便得到整个含噪数据的边缘分布
\begin{equation}
 p_t(x)=\int q_t(x\mid x_0)\,\pdata(x_0)\,\dd x_0,
 \label{eq:c-mixture}
\end{equation}
这里为便于书写，暂且假设 $\pdata$ 有密度；对离散或奇异分布，应把积分写为对概率测度的积分，含义相同。加噪改变的是整群样本的分布，而不是只对一张图片做一个视觉滤镜。

如果直接输出平均图像，不同对象会叠在同一张图里；高斯加噪则使每张原图产生不同的随机观测。前者把不确定性压成一个均值，后者用分布表达不确定性。图像最终看起来像噪声，是随机扰动遮蔽了结构，而不是所有原图都收缩成同一个灰点。

\subsection{高斯平滑为什么有帮助}

为看清加噪的作用，暂设 $a(t)=1$、$b(t)=\sigma$。此时 $p_\sigma$ 是原分布与高斯核的卷积。原来尖锐、分离或者没有普通密度的结构，经过正方差的高斯扰动后，会形成一个平滑、处处为正的密度。这让“在某个观测附近，概率怎样变化”有了良好定义。

小噪声保留细部差异，每一簇数据附近仍有自己的局部结构；大噪声使相邻簇重叠，细节逐渐难以辨认。对于两个分开的数据簇，先研究较宽、较平滑的概率地形，再逐步增加分辨率，通常比从远离数据的任意点直接寻找细小簇更有希望。这个多尺度安排正是理解扩散的重要入口，但它不是对有限模型必然成功的证明。

\begin{figure}[htbp]
\centering
\begin{minipage}{0.48\linewidth}
\centering
\begin{tikzpicture}
\begin{axis}[width=\linewidth,height=4.5cm,domain=-4.5:4.5,samples=120,
 xmin=-4.5,xmax=4.5,ymin=0,ymax=0.55,axis lines=left,
 xlabel={$x$},ylabel={$p_\sigma(x)$},title={较小噪声 $\sigma=0.4$},
 tick label style={font=\footnotesize},label style={font=\small},
 title style={font=\small}]
\addplot[blue!70!black,line width=1.2pt]
 {0.5/(sqrt(2*pi)*0.4)*(exp(-(x+2)^2/(2*0.4^2))+exp(-(x-2)^2/(2*0.4^2)))};
\end{axis}
\end{tikzpicture}
\end{minipage}\hfill
\begin{minipage}{0.48\linewidth}
\centering
\begin{tikzpicture}
\begin{axis}[width=\linewidth,height=4.5cm,domain=-6:6,samples=120,
 xmin=-6,xmax=6,ymin=0,ymax=0.55,axis lines=left,
 xlabel={$x$},ylabel={$p_\sigma(x)$},title={较大噪声 $\sigma=2.2$},
 tick label style={font=\footnotesize},label style={font=\small},
 title style={font=\small}]
\addplot[orange!85!black,line width=1.2pt]
 {0.5/(sqrt(2*pi)*2.2)*(exp(-(x+2)^2/(2*2.2^2))+exp(-(x-2)^2/(2*2.2^2)))};
\end{axis}
\end{tikzpicture}
\end{minipage}
\caption{对等概率的两个点 $-2$ 与 $2$ 添加不同强度的高斯噪声。右图的密度更平滑，两个来源更难区分。本图改变了噪声标准差，未进行信号缩放；因此它说明平滑作用，而不是完整的方差保持扩散过程。}
\label{fig:c-smoothing}
\end{figure}

回到式\eqref{eq:c-noising}。当 $a(t)$ 足够小、$b(t)$ 接近 $1$ 时，$X_t$ 的分布接近标准高斯。这里的“接近”有必要保留：有限时间的 DDPM 通常只有近似高斯的末端分布。若 $a(1)=0$ 且 $b(1)=1$，则末端才精确地独立于 $X_0$。生成前为初始随机状态指定的分布称为生成先验；选择一个容易采样、又尽量接近末端的先验，是生成链的第一环。

从矩的角度看，
\[
 \E X_t=a(t)\E X_0,\qquad
 \operatorname{Cov}(X_t)
 =a(t)^2\operatorname{Cov}(X_0)+b(t)^2I_d,
\]
假设相应矩存在。信号系数趋于零时，数据的均值与协方差影响逐渐减弱。仅仅矩接近还不等于分布完全接近，但这能解释为什么标准高斯是一个合理的终点设计。

\subsection{失去信息之后，反向怎么可能？}

如果加噪逐渐销毁信息，从纯噪声出发怎么可能还原图片？困惑来自“还原”的两个含义。恢复加噪前那一张特定原图，需要观测中保留足够信息；产生一个符合原分布的新样本，却只需要知道怎样对所有可能的原图进行选择。扩散生成通常做的是后者。

设观察到一张模糊而含噪的图。它可能来自窗台上的猫，也可能来自沙发上的猫。一个生成模型无须知道历史上到底是哪张图片被扰乱，而要知道哪些解释合理、各自有多少概率。每次反向采样可选出一个不同解释，重复运行后再现总体规律。信息没有从噪声中凭空出现；关于图像结构的信息来自训练数据和已经学到的参数，随机性负责选择具体结果。

这也解释了为何加噪过程可以是已知的，而反向过程需要学习。给原图加多少噪声，我们可以控制；给定噪声观测，哪种结构更合理，则需要了解数据分布。若把 $\pdata$ 换成手写数字、医学切片或者乐谱，正向高斯规则仍可相同，反向规律却完全不同。

在贝叶斯语言中，尚未考虑当前观测时的规律称为先验；结合观测之后，对未知原图的条件分布称为后验。这里 $\pdata(x_0)$ 提供原图先验，$q_t(x_t\mid x_0)$ 描述观测如何产生，二者通过贝叶斯公式得到后验 $p(x_0\mid x_t)$。后文讨论的后验均值和后验采样，都是针对这份经过更新的概率规律。它与生成链末端的噪声先验承担不同角色。

\section{去噪为什么会携带分布的信息？}
\label{sec:c-denoising}

\subsection{训练标签可知，不代表推理时答案可知}

训练时，我们先取一张干净图 $x_0$，再自行产生 $\eps$，计算 $x_t=a(t)x_0+b(t)\eps$。因此 $x_0$ 和 $\eps$ 都是现成标签。可用网络 $D_\theta(x_t,t)$ 预测干净图，也可用 $\eps_\theta(x_t,t)$ 预测噪声。表面上看，这只是一个普通回归任务。

但是，推理时只给网络 $x_t$ 和 $t$，不给它产生这张观测的 $x_0$ 和 $\eps$。同一个 $x_t$ 可以由不同的原图与噪声组合得到，所以网络不可能一般性地精确识别这一次所加的噪声。平方误差下，它学到的理想函数是条件均值：
\[
 D^*(x,t)=\E[X_0\mid X_t=x],\qquad
 \eps^*(x,t)=\E[\eps\mid X_t=x].
\]
这不是失败，而是回归目标的正确统计含义。给定二者之一，在 $a(t),b(t)$ 非零时，可以换算另一个：
\[
 D^*(x,t)=\frac{x-b(t)\eps^*(x,t)}{a(t)}.
\]

条件均值依赖 $\pdata$。若训练数据只有水平线，去噪器会倾向解释出水平线；若数据只有圆圈，同样的模糊输入会得到不同预测。因此回归任务并非只在学一种通用滤波器，它还在学习“这种观测通常由哪些数据产生”。

\subsection{一个能算清楚的两点世界}

继续使用 $X_0\in\{-h,h\}$、各占一半的例子，设
\[
 Y=X_0+\sigma\eps,\qquad \eps\sim\N(0,1),\quad h>0.
\]
其中 $X_0$ 与 $\eps$ 独立。观察 $Y=y$ 后，对两个来源的相对支持为
\[
 \frac{\mathbb P(X_0=h\mid Y=y)}
      {\mathbb P(X_0=-h\mid Y=y)}
 =\exp\!\left(\frac{2hy}{\sigma^2}\right).
\]
因此最优平方误差去噪器为
\begin{equation}
 D^*(y)=h\tanh\!\left(\frac{hy}{\sigma^2}\right).
 \label{eq:c-two-point-denoiser}
\end{equation}
正的观测支持右边端点，负的观测支持左边端点；证据越强，预测越接近对应端点。当 $\sigma$ 较大或 $y$ 接近零时，两种解释难分，均值就停留在它们之间。

这个例子同时说明了去噪器的能力和局限。$D^*(0)=0$，而原始数据从不取 $0$。即使一个去噪器统计上完全正确，一次输出的条件均值也不必是数据分布中的合法样本。高噪声时“预测的干净图很模糊”，并不直接说明模型没学好；可能只是后验不确定性很大。问题是如何利用这个预测构造正确的反向过程，而不是期待一个均值估计器一次完成全部生成。

对相同含噪观测，后验采样可以返回 $-h$ 或 $h$，后验均值则返回它们的加权平均。两种运算回答的是不同问题。扩散采样会通过一系列更新逐步降低不确定性，不能简单等同于对每一步都“直接输出预测原图”。具体更新式由时间路径、反向分布和采样器共同决定。

\subsection{去噪残差是一种概率方向}

令 $p_\sigma(y)$ 是 $Y=X_0+\sigma\eps$ 的密度。一个关键恒等式是
\begin{equation}
 \grad_y\log p_\sigma(y)
 =\frac{D^*(y)-y}{\sigma^2}.
 \label{eq:c-tweedie-intuition}
\end{equation}
左边称为\term{分数函数}{gl:score}，通常简称 score。这个词与“给图片打一个分数”无关：它是对数密度关于观测坐标的梯度，是一个与图片同维度的向量。式\eqref{eq:c-tweedie-intuition}告诉我们，去噪器建议的修正方向包含了含噪分布的局部信息。去噪与 score matching 的联系是扩散的重要思想来源。\cite{c_vincent2011connection}

读者熟悉梯度就可以理解其局部含义。在一个很小的位移 $\Delta y$ 下，
\[
 \log p_\sigma(y+\Delta y)-\log p_\sigma(y)
 \approx \ip{\grad\log p_\sigma(y)}{\Delta y}.
\]
score 指向密度局部增加最快的方向，同时给出变化的大小。它告诉运动机制在每个位置应该怎样响应，而不需要直接输出一个巨大的密度函数。完整恒等式及不同噪声系数下的版本将在第\ref{sec:score}节推导。

仍以两点世界为例，代入式\eqref{eq:c-two-point-denoiser}就有
\[
 \score_\sigma(y)
 =\frac{h\tanh(hy/\sigma^2)-y}{\sigma^2}.
\]
小噪声时，在两个端点附近，箭头把偏离簇中心的点拉回；在两簇之间，它又根据观测位置区分两种可能性。噪声尺度改变，整个向量场也改变。网络接收 $t$，就是为了知道现在应当使用哪一张“概率地图”。

\begin{cautionbox}{沿梯度上山不是分布采样}
若只反复沿 score 做梯度上升，样本通常会向密度极大值集中，不能保证保持正确的概率比例。生成需要一个使整群样本按照目标分布演化的机制。随机项、扩散系数、时间相关的漂移，以及确定性流中的体积变化，都会参与这个机制；它不是普通的图片质量优化。
\end{cautionbox}

这里也能解释“梯度不知道归一化常数，为何还足够”。若 $p(x)=\widetilde p(x)/C$，其中 $C$ 与 $x$ 无关，则 $\grad\log p(x)=\grad\log\widetilde p(x)$。局部梯度不必计算这个常数。不过，只知道梯度并不等于可以随意选择采样器；在断开的支持上，相对概率质量也有额外问题。正方差的高斯扰动把分布平滑到整个空间，有助于避开这些几何障碍。

\section{局部回归怎样组织成全局生成？}
\label{sec:c-global}

\subsection{很多简单任务，共用一个网络}

训练时，我们可以随机选噪声级别 $t$，直接由式\eqref{eq:c-noising}获得观测。无需每次都真的从 $0$ 一步一步加噪到 $t$。一个时间条件网络在大量样本和噪声级别上学习，得到一组相互关联的预测规则。小噪声任务偏向恢复细部，大噪声任务更依赖总体结构，但这只是常见现象，不是每个采样步都具有固定语义职责的定理。

许多时间任务共享参数，既节省资源，也使模型能利用尺度之间的共同规律。可以把它想象成一位修复师：轻微损伤时，看清的细节较多，需要微调；严重损伤时，必须更多依靠对对象的总体知识。但真正的网络未必像人一样先认出一个离散语义再动笔，它执行的是连续函数计算。

最简单的噪声预测训练目标具有形式
\[
 \E_{t,X_0,\eps}
 \left[\norm{\eps-\eps_\theta(a(t)X_0+b(t)\eps,t)}^2\right].
\]
这条式子很好写，也容易构造训练数据，但它的作用不能只凭直觉判断。为何这种目标与概率建模有关，不同时间是否应有不同权重，以及怎样换算成反向均值，都是第\ref{ch:math}章要完整回答的问题。形式简单，不代表背后的对象简单。

\subsection{训练和采样走在不同方向上}

\begin{table}[htbp]
\centering
\caption{把训练与采样分开理解。二者复用同一个网络，但输入来源、已知信息与任务方向不同。}
\label{tab:c-training-sampling}
\small
\begin{tabularx}{\linewidth}{@{}p{0.17\linewidth}XX@{}}
\toprule
 & 训练 & 生成采样\\
\midrule
起点 & 真实数据 $x_0$ 与自行抽取的噪声 & 一份从简单先验抽取的随机状态\\
已知信息 & 原图、所加噪声和噪声级别 & 当前状态、时间和可能的外部条件\\
网络任务 & 学习条件均值、score 或相关预测量 & 为采样器提供当前状态的更新信息\\
时间使用 & 通常随机选择时间，无须依次遍历 & 按设定的时间序列或连续积分推进\\
最终目的 & 调整参数，使许多局部预测准确 & 将初始样本群搬运为目标分布\\
\bottomrule
\end{tabularx}
\end{table}

若真实反向转移分布已知，并且从正确的末端分布出发，逐步使用这些转移，就会返回正确的数据分布。这本质上是联合分布的条件概率分解，而不是“噪声会自动变成图片”的神奇规律。学习模型用有限网络近似反向规律，采样器再用有限步数近似所需过程。训练误差、末端近似和求解误差都会改变最终输出。

因此扩散有效的直觉可以分成三层。第一，正向过程人为规定，训练样本容易构造。第二，邻近噪声级别之间的变化比较小，可以把困难生成拆成较温和的预测。第三，概率论给出这些局部规则应怎样组织成全局过程。真正把三层接起来的是数学；直觉负责让公式的用途变得可见。

请特别注意“邻近”这个词。跨越很大噪声间隔时，反向条件分布可能高度多峰，一个简单高斯未必表达得好。小时间步下的局部近似更可信，但代价是采样步骤增多。这正是后来快速采样、高阶求解、蒸馏和路径设计反复面对的张力。扩散不是免费消除复杂性，而是把复杂性分配到一条有结构的路径上。

\subsection{随机性不是误差，它承担选择的工作}

同一句提示词，多次运行模型得到不同构图，通常是我们想要的行为。未被提示词规定的姿势、背景细节、纹理等，都需要在合理范围内随机选择。初始噪声提供这种随机性；在随机采样过程里，中途噪声也参与选择。在确定性采样过程里，初始随机向量已经足以产生结果差异。

随机种子只是产生伪随机序列的一个索引，不是模型内部的“图片编号”，也不是可解释的语义编码。改变种子会改变采样轨迹，不能保证只改变一个指定对象。固定种子有助于比较参数修改的效果，但这种控制依赖于相同的模型、时间序列和实现；换了采样设置，同一个种子也未必对应同一构图。

另外，反向随机过程不必逐样本重走当初那条正向加噪路径。正确性通常关心各时刻和最终的分布，而不是每个噪声粒子能否找回“自己的原图”。这一区分将把 DDPM、DDIM 和确定性概率流统一到同一个理解框架中。

\section{隐空间：哪些东西被隐藏了？}
\label{sec:c-latent}

\subsection{同一个词，三种常见用法}

\term{隐变量或潜变量}{gl:latent}是模型中未被直接观测、却参与描述或产生数据的变量。“隐空间”则是这些变量所处的空间。它并不专指一种网络，也不意味着人类看不懂的地方一定叫隐空间。讨论扩散时，至少应区分三种用法。

第一种是生成器的随机输入 $Z$。例如 GAN 用它决定输出的变化；$Z$ 未作为训练图像的直接标签被观测，因此可称为隐变量。它的维数、分布和语义由设计与训练共同决定。高斯输入不保证各坐标变成独立语义因素。

第二种是概率模型中解释观测的未观测变量。VAE 用 $z$ 参与 $p_\theta(x\mid z)$；DDPM 的 $x_1,\ldots,x_K$ 对观测 $x_0$ 而言也是一串辅助隐变量。于是，像素空间扩散仍可被称为隐变量模型。这里“隐”说的是概率角色，不是在说它已经压缩了图像。

第三种是潜空间扩散中的压缩表示。一个编码器 $E$ 把图片变成 $z_0=E(x_0)$，扩散在这个表示空间里学习；一个解码器 $D$ 把生成的 $z_0$ 转回图片。它把表示学习与分布学习拆成两个环节，是 latent diffusion 的主要安排。\cite{c_rombach2022ldm}

这三种用法可以同时出现在同一个系统里。最初随机抽取的 $z_K$ 是随机状态；去噪中的 $z_k$ 是轨迹上的隐变量；最终 $z_0$ 是需要解码的压缩表示。符号常都写作 $z$，但它们的统计角色和时间位置不同。阅读论文时先问“这个 $z$ 是编码结果、先验样本还是含噪状态”，会减少许多误解。

\subsection{压缩表示为什么值得做}

很多像素变化对理解画面并不同等重要。相邻像素的细微差别、某些纹理和压缩噪声，可能花费大量维数，却不改变对象与构图。如果编码器能把重要结构放进较小的空间，扩散网络便可以在较少的位置上运算，最后由解码器补回可见细节。

作一个维数比较：一张 $512\times512\times3$ 图片有 $786\,432$ 个标量；若表示为 $64\times64\times4$ 的潜张量，则有 $16\,384$ 个标量，数量为原来的 $1/48$。这说明表示大小可以显著减少，但实际运行时间不必减少 $48$ 倍，因为通道宽度、网络深度、注意力、编码解码成本和硬件行为都在影响计算。

这里的 $64\times64\times4$ 只是一个示例，不是所有潜空间扩散的固定格式。潜变量也可以有不同空间压缩率、通道数、离散编码或时间维度。对视频，压缩时间轴可以降低开销，同时也可能损失快速运动细节。不同表示的选择，决定了后面的生成器容易保留什么、难以恢复什么。

\begin{figure}[htbp]
\centering
\begin{tikzpicture}[x=1cm,y=1cm]
 \node[draw,rounded corners,fill=blue!6,minimum width=2.3cm,minimum height=1.0cm,align=center] (img) at (0,0) {像素图像\\$x_0$};
 \node[draw,rounded corners,fill=blue!6,minimum width=2.3cm,minimum height=1.0cm,align=center] (lat) at (3.6,0) {压缩表示\\$z_0=E(x_0)$};
 \node[draw,rounded corners,fill=orange!8,minimum width=2.3cm,minimum height=1.0cm,align=center] (noise) at (7.2,0) {含噪表示\\$z_k$};
 \draw[->,thick] (img) -- node[above,font=\small] {编码} (lat);
 \draw[->,thick] (lat) -- node[above,font=\small] {加噪} (noise);
 \draw[->,thick] ([yshift=-0.25cm]noise.west) to[bend left=28] node[below,font=\small] {学习反向生成} ([yshift=-0.25cm]lat.east);
 \node[draw,rounded corners,fill=green!7,minimum width=2.3cm,minimum height=1cm,align=center] (out) at (3.6,-2.0) {输出图像\\$\widehat x_0=D(\widehat z_0)$};
 \draw[->,thick] (lat) -- node[right,font=\small] {解码} (out);
 \node[font=\small,align=center] at (0,-2.0) {表示决定信息瓶颈\\扩散学习潜变量分布};
\end{tikzpicture}
\caption{潜空间扩散的两个学习层次。图中编码与解码说明表示的角色；生成时从潜噪声开始，并不先输入一张真实图片。}
\label{fig:c-latent}
\end{figure}

\subsection{潜空间不是免费的完美语义空间}

压缩意味着取舍。若 $D(E(x))$ 已经把小字、精细线条或某种医学特征损坏，后面的扩散模型就不能普遍保证把它们准确恢复。它可能生成一个看起来合理的细节，但合理并不等于与原始事实一致。潜空间扩散的输出能力受到表示的重建能力和解码器行为共同限制。

从概率角度看，扩散首先学习的是训练图像经编码后形成的潜变量分布。设 $Z_0=E(X_0)$，理想生成结果满足 $\widehat Z_0$ 与 $Z_0$ 同分布，那么解码后复现的是 $D(Z_0)$ 的分布。除非编码解码在目标任务下足够保真，它并不精确等于原始 $X_0$ 的分布。这个推送关系比“潜空间里学好了，图像就自然学好了”更准确。

同样，给每张图像的编码加一个高斯正则约束，也不意味着整批编码形成的分布精确是标准高斯。单个样本的近似后验、所有样本的聚合分布和生成先验是不同对象。我们可以专门在潜空间里训练扩散，恰恰是为了学习那个复杂的聚合分布。第\ref{ch:applications}章会进一步讨论表示模型与生成主干的合作。

\section{条件生成：把“可能”缩小到“符合要求”}
\label{sec:c-conditional}

\subsection{一句描述通常对应许多图片}

无条件模型学习 $\pdata(x)$；条件模型学习 $p_\theta(x\mid c)$。条件 $c$ 可以是类别、文字、草图、深度图、已有视频片段或测量结果。一句“窗边的一只橘猫”约束了部分内容，却没有规定相机角度、猫的动作、窗户形状以及每一根毛的状态。于是条件生成仍然是分布问题，不能一般性地压成唯一正确图像。

训练时若有配对样本 $(x_0,c)$，就可以给网络额外输入 $c$，学习
\[
 \eps_\theta(x_t,t,c)\approx \E[\eps\mid X_t=x_t,C=c].
\]
从 score 的角度看，它在估计 $\grad_x\log p_t(x\mid c)$。条件越有信息，允许的解释一般越受限制；但“信息多”不保证数据中的条件可靠，也不保证网络准确使用了每个条件。

如果提示词遗漏某种属性，模型会依据训练分布补充；如果提示词要求一个训练中很少见的组合，模型可能只完成其中一部分。这说明条件不只是用户界面的输入，还决定了概率建模的可识别性与泛化难度。词语组合能够表达的空间通常远大于有限配对数据实际覆盖的空间。

\subsection{条件是怎样进入图像网络的}

条件可以转换成特征，再通过拼接、调制或者注意力与图像表示发生交互。文字编码器会把文本转换为一组\term{词元}{gl:token}表示；交叉\term{注意力}{gl:attention}让图像位置根据当前状态读取有关的文本信息。此处重点是信息通道：概率目标说明应学习什么关系，具体结构说明怎样把条件送入计算。

文本编码器强，不意味着生成结果一定能逐字遵循指令。生成主干仍要把语言关系映射到空间关系；多个对象的属性绑定、数量、位置和细小文字都可能出错。条件表示、训练配对和网络交互方式分别影响不同环节。第\ref{ch:applications}章会从具体模型中看到这些问题怎样被处理。

另一个常见工具是 guidance，中文通常称“引导”。它在采样时加强条件的影响。第\ref{sec:guidance}节将从概率角度说明其由来，第三章再比较具体实现。现在只需记住：增强提示词符合度通常会改变输出分布，可能同时改变多样性、饱和度与细节。调高引导强度不等于无条件地提高模型的全部能力。

\subsection{图像修复也是条件概率问题}

设真实图像 $X$ 经测量算子 $A$ 得到 $Y=AX+\eta$，例如 $A$ 只保留部分像素，或执行低分辨率采样。恢复任务应研究 $p(x\mid y)$。给定缺损图像，缺失区域可能有多种合理补全；给定低分辨率图像，高频纹理也可能无法唯一确定。生成式修复正是用数据先验补充测量没有提供的信息。

对于拍照滤镜，我们可能欢迎丰富而美观的补充；对于科学测量，则要区分观测支持的结构与先验补出的结构。数学上的条件分布让这一区分可以被清楚表达：结果的某些部分被数据约束，另一些部分仍然存在不确定性。理解扩散的能力，也需要理解它对不确定性的使用方式。

\section{把扩散放回生成模型的全景}
\label{sec:c-families}

扩散不是第一个学习图像分布的模型，也不是把所有其他方法简单替换掉的一条单线进化。不同家族选择不同的概率表达、学习信号和采样计算；许多现代系统会混合多个家族。下面按“如何组织分布”理解它们，而不是只记住几个缩写。

\subsection{自回归：把联合概率拆成条件概率}

\term{自回归模型}{gl:ar}利用概率链式法则，
\[
 p_\theta(x_1,\ldots,x_d)
 =\prod_{i=1}^{d}p_\theta(x_i\mid x_1,\ldots,x_{i-1}).
\]
先决定顺序，再预测下一个位置。对图像，原始 PixelRNN 系列采用像素级的条件建模；对于现代视觉系统，序列元素也可以是压缩后的离散词元。\cite{c_oord2016pixelrnn}这种方式使似然分解直接，训练可以通过已知前缀提供监督，采样则依赖顺序生成。

它的关键选择是顺序。图片本来是二维整体，却要被排成一个序列；之前生成的内容成为之后生成的条件。扩散则通常在每一步同时更新所有位置，通过多次整体处理协调它们。这里的“同时”指状态更新，不意味着每个位置互相独立。

不应简单地说自回归一定慢、扩散一定快。逐像素采样可能有很长序列，压缩词元会缩短它；扩散每一步能并行处理位置，却常需要多次调用网络。实际速度由序列长度、采样步数、模型大小和硬件共同决定。概率分解的不同，让二者的计算瓶颈也不同。

\subsection{VAE：通过潜变量解释数据}

\term{变分自编码器}{gl:vae}首先规定
\[
 p_\theta(x)=\int p(z)p_\theta(x\mid z)\,\dd z,
\]
然后用编码分布 $q_\phi(z\mid x)$ 近似潜变量后验。这样既有生成路径，也有从观测推断潜变量的路径。由于积分常不可直接计算，训练使用\term{证据下界}{gl:elbo}；重参数化使其中的重要期望可以通过梯度进行优化。\cite{c_kingma2014vae}

普通自编码器只学编码和重建，并不自动保证“随机抽一个潜变量就能得到合理图片”。VAE 增加了概率模型和分布约束，但这些约束同样存在表达能力与重建精度的取舍。某些 VAE 的模糊输出与简单解码分布及平方误差有关，不能把“所有 VAE 都模糊”当成家族定义。

DDPM 也能从含有许多辅助隐变量的模型出发写出变分下界。因此扩散与 VAE 不是完全无关的两个数学世界。不过，标准 VAE 常用较直接的潜变量到图像解码，扩散把生成分配到许多噪声尺度；潜空间扩散又可以用一个 VAE 类表示模型承接图像与潜变量。关系在概率建模层面与系统组成层面各有一层。

\subsection{GAN：通过对抗反馈学习一个采样器}

\term{生成对抗网络}{gl:gan}把随机输入经过生成器 $G_\theta$ 变成样本，再让判别器区分真实与生成样本。生成器根据判别器提供的反馈更新，从而学习一个能产生目标分布样本的映射。\cite{c_goodfellow2014gan}它通常无须计算输出的归一化密度，采样也可以只需一次主生成器调用。

这种安排能让学习信号强调真假可辨的细部，因此在图像生成史上取得了重要成果；同时，两个网络互相影响会造成优化上的困难，生成器还可能忽略一些数据模式。所谓“模式”是某类具有独特内容或结构的可能性，例如一类姿态或一个数字类别，而不只是纹理风格。模式缺失就是输出分布没有充分覆盖它们。

扩散把训练组织成已知扰动下的预测，通常不需要判别器参与核心目标。这有助于避免某些对抗训练难题，但不保证不会发生覆盖不足、训练失衡或记忆。后来的快速扩散模型也可以结合对抗损失。因此“使用去噪”与“使用判别器”并非互斥的系统属性。

\subsection{归一化流：通过可逆变换管理密度}

\term{归一化流}{gl:flow}取可逆映射 $x=f_\theta(z)$，用变量替换公式写出
\[
 \log p_\theta(x)
 =\log p_Z(f_\theta^{-1}(x))
 +\log\left|\det J_{f_\theta^{-1}}(x)\right|.
\]
这里 $J$ 是\term{雅可比矩阵}{gl:jacobian}。它描述变换怎样拉伸或压缩体积，因此密度需要随体积改变。Real NVP 等模型通过特殊结构使可逆计算和行列式计算可行。\cite{c_dinh2017realnvp}

归一化流把采样和密度计算联系得十分直接，但计算方便的可逆结构也约束了设计。与之相比，常见去噪网络不必自身可逆；整个生成过程的数学规律负责组织分布变换。连续的确定性流与概率流 ODE 又会把这些思路重新联系起来。

还有一个术语陷阱：“normalizing flow”中的 flow，与“flow matching”不是同一个方法名。前者强调用可逆变换和体积变化表达密度；后者强调用回归目标学习一个连续速度场。二者可以相交，例如由满足适当条件的速度场积分得到连续归一化流，但不能仅看到 flow 就假设模型使用同一种目标或同一种采样器。

\subsection{能量模型：先描述相对偏好，再想办法采样}

能量模型写成
\[
 p_\theta(x)=\frac{\exp(-E_\theta(x))}{Z_\theta},\qquad
 Z_\theta=\int \exp(-E_\theta(x))\,\dd x.
\]
低能量表示相对偏好；配分函数 $Z_\theta$ 使密度归一化。一个任意的能量函数未必拥有有限配分函数，所以数学上还需要可积性。对高维神经能量模型，归一化以及训练中的采样往往很难。

能量梯度与 score 有直接联系：$\grad_x\log p_\theta(x)=-\grad_xE_\theta(x)$。模型可以通过这类局部方向参与随机采样。现代神经能量模型提供了这种思路的具体实例。\cite{c_du2019ebm}扩散 score 模型则对一系列噪声分布分别学习方向，通常不需要在每次训练更新中先运行完整的生成链。

但 score 网络也不一定显式表示为某个标量能量的精确梯度。理想目标是一个梯度场，有限网络近似后可能不完全具有对应的保守场结构。这不会阻止我们把它放入采样器，却提醒我们：一个有数学来源的目标，与有限学习结果满足全部数学性质，是两个不同层次。

\begin{table}[htbp]
\centering
\caption{生成模型家族的核心安排。比较针对典型形式，实际系统可混合多个方法。}
\label{tab:c-family-overview}
\small
\begin{tabularx}{\linewidth}{@{}p{0.13\linewidth}XXX@{}}
\toprule
家族 & 分布如何表达 & 主要学习信号 & 典型采样计算\\
\midrule
自回归 & 按顺序分解条件概率 & 已知前缀下的似然 & 依次产生序列元素\\
VAE & 潜变量与条件解码分布 & 变分下界与重建 & 抽潜变量，再解码\\
GAN & 随机输入的生成映射 & 真假判别的反馈 & 执行生成映射\\
归一化流 & 可逆变换与体积因子 & 可计算的对数密度 & 从简单分布经可逆映射\\
能量模型 & 能量与归一化常数 & 似然或替代目标 & 常需迭代随机采样\\
扩散 / score & 多噪声级别的反向规律 & 去噪、score 与变分目标 & 多步随机或确定性更新\\
\bottomrule
\end{tabularx}
\end{table}

\section{建立四层坐标，避免把名字混成一团}
\label{sec:c-four-layers}

读到新的生成模型时，建议先把它放到四层坐标上。第一层是表示：状态是像素、压缩张量还是离散词元？第二层是概率路径：数据怎样与简单分布联系起来，时间中间是什么分布？第三层是学习目标：网络预测噪声、干净数据、score 还是速度？第四层是网络与求解：怎样表达这个函数，怎样按有限步数生成？

例如，“潜空间”“DiT”“噪声预测”“DDIM”分别主要回答上述不同层的问题。它们不一定相互竞争，反而可以组合成一个系统。一个 Transformer 主干可以做噪声预测，也可以做速度回归；一个潜表示可以服务扩散，也可以服务自回归生成；换一个采样器通常不等于重新学习全部表示。

这里还需把“结构”一词说清楚。U-Net 和 DiT 主要指神经网络结构；DDPM 指一组概率建模、训练和采样安排；DDIM 主要提供与训练边缘相联系的另一类生成轨迹；flow matching 是学习连续速度场的一类方法。第三章将沿数学思想梳理它们的发展关系，而不会把这些名字放在同一张架构排名表上。

\begin{keybox}{进入第二章前，先保留这五句话}
生成的对象是分布；随机性选择具体样本。加噪给复杂分布建立了一串较平滑的中间分布。去噪的条件均值可以转换为局部概率方向。局部方向需要正确的时间机制才能产生全局采样。表示、概率路径、预测目标和网络结构是相互关联但可分别讨论的层次。
\end{keybox}

我们现在已经获得扩散可能有效的直觉，却还欠着严格答案：离散反向均值为何可以用噪声预测来表示？训练目标为何来自变分下界？score 怎样进入反向 SDE？概率流 ODE 如何让随机与确定性生成相遇？速度场又如何统一描述概率搬运？第\ref{ch:math}章将把这些欠下的推导逐一补齐。

\section{练习与解答思路}
\label{sec:c-exercises}

\begin{exercisebox}{练习一：平方误差为何偏爱平均答案？}
设 $X$ 是二阶矩有限的随机向量，模型必须输出固定向量 $m$。证明 $\E\norm{X-m}^2$ 在 $m=\E X$ 时最小。再解释：为什么这个结论不意味着生成模型应当总输出均值？
\end{exercisebox}

\begin{exercisebox}{练习二：一个有偏的两点分布}
设 $\mathbb P(X_0=h)=\pi$、$\mathbb P(X_0=-h)=1-\pi$，其中 $0<\pi<1$，并观察 $Y=X_0+\sigma\eps$。写出后验对数几率和最优去噪器。若 $\pi>1/2$，在 $Y=0$ 时去噪结果是否仍为零？这说明了什么？
\end{exercisebox}

\begin{exercisebox}{练习三：确定性函数能产生随机分布吗？}
设 $U$ 在 $(0,1)$ 上均匀分布，$X=\mathbf 1_{\{U>\frac34\}}$。求 $X$ 的分布。随后设 $G$ 是固定网络，讨论“采样过程没有中途噪声，所以不可能产生多样性”这句话哪里有误。
\end{exercisebox}

\begin{exercisebox}{练习四：理想潜空间生成能保证什么？}
设编码器为 $E$、确定性解码器为 $D$，潜空间生成器精确复现 $Z=E(X)$ 的分布。证明输出分布等于 $D(E(X))$ 的分布。给出一个简单例子，说明它不一定等于 $X$ 的分布；再给出一个充分条件，使二者相等。
\end{exercisebox}

\begin{exercisebox}{练习五：把一个系统放到四层坐标上}
某图像系统先将图片编码为连续潜张量，以高斯扰动构造训练数据，用 Transformer 预测噪声，再用确定性多步更新采样。分别指出它的表示、概率路径、学习目标与网络/求解选择。说明“它是 Transformer，所以不属于扩散”为什么不成立。
\end{exercisebox}

\subsection*{解答思路}

\paragraph{练习一。}
令 $\mu=\E X$，展开得到
\[
 \E\norm{X-m}^2
 =\E\norm{X-\mu}^2+\norm{\mu-m}^2,
\]
交叉项因 $\E(X-\mu)=0$ 而消失。均值是给定限制下的最优估计，但生成要求产生符合整个分布的随机结果，而不是在固定向量集合里挑选一个最优代表。对多峰分布，均值甚至可能落在低概率区域。

\paragraph{练习二。}
贝叶斯公式给出后验对数几率
\[
 \ell(y)=\log\frac{\pi}{1-\pi}+\frac{2hy}{\sigma^2}.
\]
设后验正端点概率为 $r(y)=1/(1+\exp(-\ell(y)))$，则
$D^*(y)=h(2r(y)-1)=h\tanh(\ell(y)/2)$。
当 $y=0$ 时，$D^*(0)=h(2\pi-1)$。观测无明显倾向时，数据先验仍会参与估计；去噪器确实携带分布的概率权重，而不只是做几何投影。

\paragraph{练习三。}
$\mathbb P(X=1)=1/4$，$\mathbb P(X=0)=3/4$。映射是确定的，但输入是随机的。相同道理适用于 $G(Z)$ 和从随机初值出发的确定性积分。多样性的质量取决于映射是否把输入概率正确组织起来，不取决于是否每一步都掷硬币。

\paragraph{练习四。}
若 $\widehat Z$ 与 $E(X)$ 同分布，则对任意可测集合 $A$，
\[
 \mathbb P(D(\widehat Z)\in A)
 =\mathbb P(E(X)\in D^{-1}(A))
 =\mathbb P(D(E(X))\in A).
\]
若 $X$ 等概率取 $-1,1$，$E$ 把两者都编码为 $0$，且 $D(0)=0$，输出就只取 $0$。一个充分条件是 $D(E(X))=X$ 几乎处处成立；“重建视觉上相近”只是近似条件，不能推出分布精确相等。

\paragraph{练习五。}
表示是连续压缩潜张量；训练路径是所指定的高斯扰动路径，还需知道具体系数才能进一步分类；预测目标是噪声；网络是 Transformer，求解是确定性多步更新。主干结构决定函数如何计算，扩散的概率安排决定它为什么这样训练和采样，二者属于不同层次。

\chaptertakeaway{本章把“生成一张图”改写为“从一个学到的分布取样”。图像的高维结构解释了学习困难，平滑与多尺度解释了加噪的作用，条件回归与 score 的联系解释了去噪为何携带概率信息。我们同时保留了三个边界：条件均值不是后验样本，密度上升不是正确采样，潜空间学得好也受编码解码的信息瓶颈限制。带着这些区分进入第二章，公式就会成为解决明确问题的工具。}
